\section{Solution sets of linear systems}
A homogeneous system has right side zero: $Ax=0$. It is always
consistent, because $x=0$ works. A nontrivial solution means any solution
other than this entire zero vector; some individual coordinates may still
be zero. A homogeneous system has a nontrivial solution exactly when
there is a free variable.

To write a solution set parametrically, solve for pivot variables in terms
of the free variables, then separate the resulting vector into its constant
part and its parameter-dependent parts. If $p$ is one solution of $Ax=b$,
all solutions are $p+v$, where $Av=0$. Thus a consistent nonhomogeneous
solution set is a translation of the homogeneous solution set.

\practice{1.5.P1}{Intersect two planes by solving their simultaneous equations}
The two equations must hold at the same time. Starting with
$x_1+4x_2-5x_3=0$ and $2x_1-x_2+8x_3=9$, subtract twice
the first from the second to get $-9x_2+18x_3=9$.
Thus $x_2=-1+2x_3$ and the first equation gives $x_1=4-3x_3$.
Write $x_3=t$. Then
\[x=(4,-1,0)^T+t(-3,2,1)^T,\qquad t\in\R.\]
Substitution gives $4-4+(-3+8-5)t=0$ in the first equation
and $8+1+(-6-2+8)t=9$ in the second. The intersection is
therefore the displayed line, not an empty set or an entire plane.

\practice{1.5.P2}{A translated plane}
Solve $10x_1-3x_2-2x_3=7$ for its pivot variable:
$x_1=7/10+(3/10)x_2+(1/5)x_3$. With $x_2=s,x_3=t$ this gives
\[x=(7/10,0,0)^T+s(3/10,1,0)^T+t(1/5,0,1)^T.\]
The two direction vectors are the same as for the homogeneous equation
in Example 2. The constant vector translates that plane; it satisfies the
new equation itself, since $10(7/10)=7$.

\exercise{1.5.1}{Find a free variable with two eliminations}
Add row 1 to row 2 and subtract twice row 1 from row 3. The
coefficient matrix becomes
\[\begin{bmatrix}2&-5&8\\0&-12&9\\0&12&-9\end{bmatrix}.\]
The last two rows are negatives. Their sum is zero, leaving only two
pivot rows for three unknowns. Hence $x_3$ is free and a nontrivial
solution exists. For an explicit check, $(-17,6,8)^T$ gives
$-34-30+64=0$, $34-42+8=0$, and $-68+12+56=0$.
\answer{Yes, the system has nontrivial solutions.}

\exercise{1.5.2}{Three pivots force the zero solution}
The operations $R_2\leftarrow R_2+2R_1$ and
$R_3\leftarrow R_3-2R_1$ give
\[\begin{bmatrix}1&-2&3\\0&-7&2\\0&0&3\end{bmatrix}.\]
All three columns have pivots. Back substitution gives $x_3=0$,
then $x_2=0$, then $x_1=0$.
\answer{No: the only solution is the trivial solution.}

\exercise{1.5.3}{Fewer equations guarantee a free variable here}
A two-row coefficient matrix can have at most two pivots. With
three unknowns, at least one variable must be free. Since the system
is homogeneous, it is consistent, and the free variable gives a
nontrivial solution. One such solution is $(-8,-4,1)^T$:
$24-16-8=0$ and $16-20+4=0$.
\answer{Yes; the pivot-count argument needs no row operations.}

\exercise{1.5.4}{A second two-equation homogeneous system}
There are three coefficient columns but only two rows. There cannot
be a pivot in every column, so a free variable exists. Homogeneity
ensures consistency, hence a nontrivial solution.
For a numerical check, $(-2,16,29)^T$ gives
$-10-48+58=0$ and $6-64+58=0$.
\answer{Yes, the system has nontrivial solutions.}

\exercise{1.5.5}{Separate the common free parameter}
The second equation is $-2$ times the first, so it adds no restriction.
The last equation gives $x_2=-x_3$. Divide the first equation by $2$
to obtain $x_1+x_2+2x_3=0$, hence $x_1=-x_3$.
Writing $x_3=t$ gives
\[x=(-t,-t,t)^T=t(-1,-1,1)^T,\qquad t\in\R.\]
Substitution into each original row gives zero.
\answer{The solution set is a line through the origin.}

\exercise{1.5.6}{Three equal coordinates}
The third equation gives $x_2=x_1$. Substituting into the first
produces $3x_1-3x_3=0$, so $x_1=x_3$. The second equation then
also holds, since $2x_3+x_3-3x_3=0$.
Thus, with $x_3=t$,
\[x=t(1,1,1)^T,\qquad t\in\R.\]
\answer{This is another line through the origin.}

\exercise{1.5.7}{Two free variables in four coordinates}
Row equivalence preserves solutions to the homogeneous equation.
The second displayed row gives $x_2=4x_3-5x_4$; substitute into the
first to obtain $x_1=-9x_3+8x_4$. Let $x_3=s$ and $x_4=t$.
\[x=s(-9,4,1,0)^T+t(8,-5,0,1)^T,\qquad s,t\in\R.\]
These expressions satisfy both displayed rows, and every solution
arises because its free coordinates may be chosen arbitrarily.

\exercise{1.5.8}{Solve the second pivot row first}
The second row gives $x_2=-2x_3+4x_4$. The first gives
$x_1=3x_2+8x_3-5x_4=2x_3+7x_4$. Hence
\[x=s(2,-2,1,0)^T+t(7,4,0,1)^T,\qquad s,t\in\R.\]
Here $s=x_3$ and $t=x_4$ are independent parameters; neither is
silently set to zero.

\exercise{1.5.9}{The third coordinate is a pivot variable}
The equations are $3x_1-6x_2+6x_3=0$ and
$-2x_1+4x_2-2x_3=0$. Twice the first plus three times the second
gives $6x_3=0$. Thus $x_3=0$ and $x_1=2x_2$.
\[x=t(2,1,0)^T,\qquad t\in\R.\]
The parameter is $x_2$, not $x_3$: the latter has been forced to zero.

\exercise{1.5.10}{A missing column gives a completely free coordinate}
The equations are $-x_1-4x_2-4x_4=0$ and
$2x_1-8x_2+8x_4=0$. Adding twice the first to the second gives
$-16x_2=0$. Hence $x_2=0$ and $x_1=-4x_4$.
The variable $x_3$ never appears, so it is independently free.
\[x=s(0,0,1,0)^T+t(-4,0,0,1)^T,\qquad s,t\in\R.\]

\exercise{1.5.11}{Three free coordinates in a six-coordinate vector}
The third nonzero row gives $x_5=4x_6$ and the second gives
$x_3=x_6$. Substitution into the first yields
\[x_1=4x_2+2x_3-3x_5+5x_6=4x_2-5x_6.\]
The free variables are $x_2=s$, $x_4=t$, and $x_6=u$. Therefore
\[x=s(4,1,0,0,0,0)^T+t(0,0,0,1,0,0)^T
+u(-5,0,1,0,4,1)^T.\]
All three parameters range independently over $\R$. The zero fourth
row imposes no additional equation.

\exercise{1.5.12}{A pivot in the last coefficient column is not inconsistency}
The third row is $x_6=0$; this is a coefficient matrix for $Ax=0$,
not an augmented matrix with a nonzero right-side pivot.
The second row gives $x_4=-4x_5+6x_6=-4x_5$.
The first then gives
$x_1=2x_2-3x_3+6x_4-5x_5=2x_2-3x_3-29x_5$.
Thus
\[x=s(2,1,0,0,0,0)^T+t(-3,0,1,0,0,0)^T
+u(-29,0,0,-4,1,0)^T,\qquad s,t,u\in\R.\]
The parameters are precisely the free coordinates $x_2,x_3,x_5$.

\exercise{1.5.13}{Read the point and direction of a line}
Put $x_3=t$ in the given coordinate description:
\[x=(5+4t,-2-7t,t)^T=(5,-2,0)^T+t(4,-7,1)^T.\]
Thus the line passes through $(5,-2,0)^T$ and has nonzero direction
$(4,-7,1)^T$. Every real $t$ is allowed.

\exercise{1.5.14}{A line can lie in four-dimensional space}
With $x_4=t$, the given formulas assemble into
\[x=(5t,3-2t,2+5t,t)^T
=(0,3,2,0)^T+t(5,-2,5,1)^T.\]
There is just one free parameter and one nonzero direction vector,
so the set is a line in $\R^4$. A line does not require its ambient
space to have only two or three dimensions.

\exercise{1.5.15}{Parallel planes from the same left side}
Solve for $x_1$ and set $x_2=s,x_3=t$. The homogeneous equation gives
\[x=s(-5,1,0)^T+t(3,0,1)^T.\]
The equation with right side $-2$ gives
\[x=(-2,0,0)^T+s(-5,1,0)^T+t(3,0,1)^T.\]
The same two independent directions occur in both descriptions.
The second plane is the first translated by $(-2,0,0)^T$.
They are distinct and disjoint: the same point cannot make the
left-hand expression equal both $0$ and $-2$.

\exercise{1.5.16}{Translate the homogeneous plane by a particular solution}
Using $x_2=s,x_3=t$, the homogeneous solutions are
\[x=s(2,1,0)^T+t(-3,0,1)^T.\]
For right side $4$ the solutions become
\[x=(4,0,0)^T+s(2,1,0)^T+t(-3,0,1)^T.\]
Both are planes with the same two directions. One passes through
zero, while the translated plane does not, since the left side at
zero is $0$, not $4$.

\exercise{1.5.17}{Translate the line from Exercise 5}
The last equation gives $x_2=-4-x_3$. Divide the first by $2$:
$x_1+x_2+2x_3=4$, so $x_1=8-x_3$. The second equation is again
$-2$ times the first and adds nothing. With $x_3=t$,
\[x=(8,-4,0)^T+t(-1,-1,1)^T,\qquad t\in\R.\]
The constant vector satisfies the original right side $(8,-16,12)^T$.
The direction agrees with Exercise 1.5.5, so the two solution lines
are parallel; the current line is translated away from the origin.

\exercise{1.5.18}{Translate the line from Exercise 6}
The third equation gives $x_2=x_1-8$. The first becomes
$3x_1-16-3x_3=5$, so $x_1=7+x_3$ and $x_2=-1+x_3$.
The second equation checks: $2(7+x_3)+(-1+x_3)-3x_3=13$.
Thus
\[x=(7,-1,0)^T+t(1,1,1)^T,\qquad t\in\R.\]
This line has the same direction as the homogeneous line in
Exercise 1.5.6 and is its translation by $(7,-1,0)^T$.

\exercise{1.5.19}{A line through a point with a specified direction}
A point on the required line is obtained by starting at $a$ and
moving any real multiple of $b$. Thus
\[x=(-2,0)^T+t(-5,3)^T,\qquad t\in\R.\]
At $t=0$ the point is $a$; differences between two points are scalar
multiples of the nonzero vector $b$, so the line is parallel to $b$.

\exercise{1.5.20}{A line through a point with a specified direction}
A point on the required line is obtained by starting at $a$ and
moving any real multiple of $b$. Thus
\[x=(3,-2)^T+t(-7,6)^T,\qquad t\in\R.\]
At $t=0$ the point is $a$; differences between two points are scalar
multiples of the nonzero vector $b$, so the line is parallel to $b$.

\exercise{1.5.21}{Use the difference of two points as direction}
Subtract the starting point from the endpoint:
$q-p=(1,4)^T$. Hence the line is
\[x=(3,-3)^T+t(1,4)^T,\qquad t\in\R.\]
The value $t=0$ gives $p$ and $t=1$ gives $q$; all displacements
are parallel to $q-p$.

\exercise{1.5.22}{Use the difference of two points as direction}
Subtract the starting point from the endpoint:
$q-p=(3,-5)^T$. Hence the line is
\[x=(-3,2)^T+t(3,-5)^T,\qquad t\in\R.\]
The value $t=0$ gives $p$ and $t=1$ gives $q$; all displacements
are parallel to $q-p$.

\exercise{1.5.23}{Five statements about homogeneous and translated systems}
\begin{enumerate}[label=(\alph*)]
\item \textbf{True.} $x=0$ always solves a homogeneous equation.
\item \textbf{False.} $Ax=0$ is an implicit condition. Parametric
vector form is an explicit description in terms of free parameters.
\item \textbf{False.} The trivial solution exists whether or not
there is a free variable. A free variable is equivalent to existence
of a \emph{nontrivial} solution.
\item \textbf{False.} $x=p+tv$ describes the line through $p$
parallel to $v$, not in general a line through $v$ parallel to $p$.
For example $p=(1,0)^T,v=(0,1)^T$ give the vertical line $x_1=1$.
\item \textbf{False as stated.} Such a representation requires $p$
to be a particular solution and $Ax=b$ to be consistent. Without
those hypotheses it need not hold. With them it is Theorem 6.
\end{enumerate}

\exercise{1.5.24}{Five further statements and their justifications}
\begin{enumerate}[label=(\alph*)]
\item \textbf{False.} A homogeneous system has the solution $0$, so
it cannot be inconsistent.
\item \textbf{False.} Nontrivial means not \emph{all} coordinates
are zero. For example $(1,0)^T$ is a nontrivial solution of
$[0\ 1]x=0$ even though its second entry is zero.
\item \textbf{True.} The displacement from $v$ to $v+p$ is
$(v+p)-v=p$, so it is in the direction of $p$ (with no movement
when $p=0$).
\item \textbf{True.} If $0$ solves $Ax=b$, then $b=A0=0$;
that is exactly the homogeneous case.
\item \textbf{True.} Consistency supplies a particular solution $p$;
Theorem 6 states the entire solution set is $p+\{v:Av=0\}$.
\end{enumerate}

\exercise{1.5.25}{Prove the translation theorem in both directions}
\goal Assume $Ap=b$. Relate every solution to the homogeneous system.
\step{(a) Add a homogeneous solution.}
If $Av_h=0$ and $w=p+v_h$, additivity of multiplication gives
$Aw=Ap+Av_h=b+0=b$. Thus every proposed vector really is a solution.
\step{(b) Recover the homogeneous part from any solution.}
If $Aw=b$, let $v_h=w-p$. Then
$Av_h=A(w-p)=Aw-Ap=b-b=0$, and $w=p+v_h$.
\answer{Both inclusions hold, so the two described solution sets
are exactly equal, not merely one contained in the other.}

\exercise{1.5.26}{The zero coefficient matrix imposes no restriction}
For the $3\times3$ zero matrix, $Ax=(0,0,0)^T$ for every input
$x\in\R^3$. Thus every vector satisfies the equation.
\[x=s(1,0,0)^T+t(0,1,0)^T+u(0,0,1)^T,\qquad s,t,u\in\R.\]
\answer{The solution set is all of $\R^3$.}

\exercise{1.5.27}{Uniqueness is equivalent to a trivial homogeneous kernel}
Because a solution exists, fix one and call it $p$. By Theorem 6,
all solutions are $p+v_h$ with $Av_h=0$.
If the only homogeneous solution is $0$, the only translated
solution is $p+0=p$, so uniqueness holds.
Conversely, if there were a nonzero homogeneous $v_h$, the vector
$p+v_h\ne p$ would give a second solution. This contradicts
uniqueness.
\answer{The equivalence follows in both directions. Consistency is
needed in order to fix $p$.}

\exercise{1.5.28}{Use the numbers of pivots, rows, and columns}
There are 3 unknowns and 3 coefficient pivots.
Consequently there are 0 free variables.
\step{(a) Nontrivial homogeneous solutions.}
Every variable is a pivot variable, so the only homogeneous
solution is zero; there is no nontrivial solution.
\step{(b) Consistency for every right side.}
Every one of the 3 rows contains a pivot. Theorem 4 therefore
guarantees consistency for every $b\in\R^{3}$.

\exercise{1.5.29}{Use the numbers of pivots, rows, and columns}
There are 4 unknowns and 3 coefficient pivots.
Consequently there are 1 free variables.
\step{(a) Nontrivial homogeneous solutions.}
Since there is at least one free variable, the homogeneous
system has nontrivial solutions.
\step{(b) Consistency for every right side.}
There are fewer pivots than rows. Theorem 4 therefore says that
some right sides have no solution.

\exercise{1.5.30}{Use the numbers of pivots, rows, and columns}
There are 5 unknowns and 2 coefficient pivots.
Consequently there are 3 free variables.
\step{(a) Nontrivial homogeneous solutions.}
Since there is at least one free variable, the homogeneous
system has nontrivial solutions.
\step{(b) Consistency for every right side.}
Every one of the 2 rows contains a pivot. Theorem 4 therefore
guarantees consistency for every $b\in\R^{2}$.

\exercise{1.5.31}{Use the numbers of pivots, rows, and columns}
There are 2 unknowns and 2 coefficient pivots.
Consequently there are 0 free variables.
\step{(a) Nontrivial homogeneous solutions.}
Every variable is a pivot variable, so the only homogeneous
solution is zero; there is no nontrivial solution.
\step{(b) Consistency for every right side.}
There are fewer pivots than rows. Theorem 4 therefore says that
some right sides have no solution.

\exercise{1.5.32}{A nonzero right side excludes the origin}
If a solution set were a plane through the origin, it would contain
$x=0$. But $A0=0\ne b$ when $b\ne0$.
\answer{No. A nonzero right side makes the origin impossible as a
solution, regardless of the dimension of the solution set.}

\exercise{1.5.33}{Arrange for row sums to vanish}
Take
\[A=\begin{bmatrix}1&-1&0\\0&1&-1\\1&0&-1\end{bmatrix}.\]
This is nonzero, and multiplying by $(1,1,1)^T$ gives row sums
$1-1+0=0$, $0+1-1=0$, and $1+0-1=0$.
\answer{The displayed matrix has the required nonzero null vector.}

\exercise{1.5.34}{Make every row orthogonal to the prescribed vector}
Take
\[A=\begin{bmatrix}1&2&0\\0&1&1\\1&0&-2\end{bmatrix}.\]
The row products with $(2,-1,1)^T$ are respectively
$2-2=0$, $-1+1=0$, and $2-2=0$.
Thus $A(2,-1,1)^T=0$, and $A$ is plainly not the zero matrix.

\exercise{1.5.35}{Cancel proportional columns by inspection}
The second column is three times the first. Therefore
$-3a_1+a_2=0$. By the column definition this means
$A(-3,1)^T=0$. The vector is nonzero, so it is nontrivial.
No elimination is needed.

\exercise{1.5.36}{Cancel a different proportionality}
The columns satisfy $2a_1+3a_2=0$, since
$2(3,-6,12)^T+3(-2,4,-8)^T=(0,0,0)^T$.
Thus $(2,3)^T$ is a nontrivial solution of $Ax=0$.

\exercise{1.5.37}{An inconsistent right side is not a translated line}
Choose
\[A=\begin{bmatrix}1&-4\\0&0\end{bmatrix}.\]
The homogeneous equation is $x_1=4x_2$, with solution line
$x=t(4,1)^T$, exactly the line through $(4,1)^T$ and zero.
Now choose $b=(0,1)^T$. The second row of $Ax=b$ says $0=1$,
so its solution set is empty, not a line parallel to the original one.
\answer{There is no contradiction: Theorem 6 assumes consistency,
and this selected right side violates that assumption.}

\exercise{1.5.38}{Scalar multiples stay in the homogeneous solution set}
The scalar-multiplication rule gives
\[A(cw)=c(Aw)=c0=0.\]
This holds for every real $c$, including $c=0$ and negative values.
\answer{Every multiple $cw$ is another homogeneous solution.}

\exercise{1.5.39}{Sums and linear combinations stay in the solution set}
Additivity first gives $A(v+w)=Av+Aw=0+0=0$.
For any scalars $c,d$, combine both multiplication rules:
\[A(cv+dw)=A(cv)+A(dw)=cAv+dAw=c0+d0=0.\]
\answer{Every such linear combination is a homogeneous solution.}

\exercise{1.5.40}{A missing right side rules out any unique solution}
Suppose there were a right-side vector $y$ for which $Ax=y$ had a
unique solution. Uniqueness would mean all
three coefficient columns have pivots. A $3\times3$ matrix would
then have a pivot in every row. Theorem 4 would make $Ax=b$
consistent for every $b\in\R^3$, contradicting the specified missing
right side.
Equivalently, the missing right side implies at most two pivots,
so at least one variable is free. Every consistent system with this
coefficient matrix consequently has infinitely many solutions.
\answer{No vector on the right can produce a unique solution.}
